% C.3bh. Full root, real_weil and frobenius proof reviews passed before registration.
\paragraph{C.3bh (Proved): integral ordinary divisor restriction in positive characteristic.}
Fix an algebraically closed field $k$ of characteristic $p>0$.
We use the precise coefficient-field scope of C.3z: $R$ is an
excellent henselian normal local $k$-algebra which is a domain,
the given coefficient field identifies its residue field with $k$,
and
\[
 N=\dim R\ge1,\qquad
 X=\operatorname{Spec}R\setminus\{\mathfrak m\},\qquad
 Y=X\cap V(f),\quad 0\ne f\in\mathfrak m,
\]
with $X$ and $Y$ regular. In particular this includes the
algebraizable strict-local scope of C.3x and C.3w.
Write $\operatorname{Pic}(X)_{\rm alg}$ for ordinary algebraic
equivalence defined by actual line-bundle families on fixed
products, as in those statements. We retain integral coefficients.
The universal dimension assertions have exactly the following answers.

\begin{center}
\begin{tabular}{llll}
\toprule
$N$ & Always injective & Always surjective & Always an isomorphism\\
\midrule
$1$ & Yes & Yes & Yes\\
$2$ & No & Yes & No\\
$3$ & No & No & No\\
$\ge4$ & No & Yes & No\\
\bottomrule
\end{tabular}
\end{center}

Every negative entry is witnessed over the algebraic closure of the
prime field and hence over the fixed field $k$. The assertion is a
universal statement about dimensions; the two negative entries for
$N=3$ need not be witnessed by the same pair. In equivalent form,
universal injectivity holds exactly for $N=1$, universal surjectivity
exactly for $N\le2$ or $N\ge4$, and universal isomorphy exactly for
$N=1$.

\begin{proof}
For $N=1$, the normal noetherian local domain is a discrete valuation
ring. Its puncture is a field spectrum and $Y$ is empty, so both
Picard groups vanish. For $N=2$, $Y$ is a regular zero-dimensional
noetherian scheme, a finite disjoint union of spectra of fields.
Its Picard group vanishes. These observations prove every positive
entry in the first two rows. For $N\ge4$, C.3z proves integral
ordinary surjectivity in the stated scope, extending the
algebraizable result C.3w. These arguments prove all the positive
surjectivity entries.

We now give the two required boundary witnesses without any
transcendental choice of a point. Use the smooth plane cubics of
C.3x,
\[
 F_p=
 \begin{cases}
 y^2z+yz^2+x^3,&p=2,\\
 y^2z-x^3+xz^2,&p\ne2,
 \end{cases}
 \qquad E=V_+(F_p),\quad O=(0:1:0),\quad P=(0:0:1).
\]
Their smoothness, including in characteristics two and three, is
proved and certified in C.3x. The line bundle
$Q=\mathcal O_E(P-O)$ is nontrivial, because the Abel map
$E\to\operatorname{Pic}^0(E)$ is an isomorphism and $P\ne O$.
For $p\ne2$, $P$ is the nonzero point of order two with $y=0$ and
$x=0$; equivalently the affine function $x$ has divisor
$2(P)-2(O)$. For $p=2$, the tangent line $y=0$ meets the cubic
only at $P$ with multiplicity three, and the affine function $y$
has divisor $3(P)-3(O)$. Thus $P$ has order three in that case.
In both cases $Q$ has nonzero order prime to $p$.

Put
\[
 R_2=\left(k[x,y,z]/(F_p)\right)_{(x,y,z)}^{\rm sh},
 \qquad X_2=\operatorname{Spec}R_2\setminus\{\mathfrak m_2\}.
\]
The actual strict-cone injection proved integrally in C.3x is
\[
 \rho^*:\operatorname{Pic}^0(E)(k)
       \lhook\joinrel\longrightarrow\operatorname{Pic}(X_2)_{\rm alg}.
\]
For clarity, its coefficient argument does not discard torsion.
On the regular cone resolution a bundle trivial on the puncture is
$\mathcal O(aE)$. Restricting to $E$ gives $Q'=\mathcal O_E(1)^{-a}$;
its degree is $-3a$, so a degree-zero $Q'$ forces $a=0$ and $Q'=0$.
The family
$\mathcal O_{E\times E}(\Delta_E-\{O\}\times E)$, pulled back
to $X_2\times E$, proves ordinary algebraic triviality using an
actual fixed smooth proper curve parameter. In particular
$\rho^*Q\ne0$, also when $k=\overline{\mathbb F}_p$.

For dimension two choose $s=x$ when $p=2$ and $s=y$ otherwise.
The exact sections are
\[
 F_2|_{x=0}=yz(y+z),\qquad
 F_p|_{y=0}=-x(x-z)(x+z)\quad(p\ne2).
\]
Each is a union of three distinct lines meeting only at the vertex.
Its puncture is regular and has trivial Picard group, as proved
in C.3x. Hence the nonzero ordinary integral class $\rho^*Q$ is
in the restriction kernel. This proves the negative injectivity
entry for $N=2$ over every fixed field $k$.

For dimension three use
\[
 R_3=k[x,y,z]_{(x,y,z)}^{\rm sh},\qquad
 X_3=\operatorname{Spec}R_3\setminus\{\mathfrak m_3\},\qquad f=F_p.
\]
The puncture of the regular local ring $R_3$ has trivial Picard
group, by the exact regular-puncture lemma already used in C.3x
\citep[Lemma 52.27.5, Tag 0F2H]{D_StacksRegularPuncture}.
Strict henselization commutes with the quotient, so $Y_3=X_3\cap
V(F_p)$ is $X_2$. The nonzero ordinary integral class $\rho^*Q$
has no source preimage. This proves the negative surjectivity
entry for $N=3$, again including $\overline{\mathbb F}_p$.

Finally C.3bg constructs a nonzero ordinary integral restriction
kernel for every $N\ge3$ and every $p$, with examples over
$\overline{\mathbb F}_p$. Its subgroup
$(\mathbb Z/p)^{\min(3,N-1)}$ proves every remaining negative
injectivity entry. Each failure of injectivity or surjectivity
also excludes universal isomorphy in that dimension. The rows
exhaust all $N\ge1$ and complete the proof.
\end{proof}

The rational table is different. Over $\overline{\mathbb F}_p$
all ordinary algebraically trivial Picard classes are torsion,
so both rational groups vanish in every dimension. Over other
algebraically closed fields the rational table is exactly C.3x,
extended to the present scope by C.3z. Neither the finite
$p$-primary kernels of C.3bg nor the finite boundary points used
here contradict those rational conclusions.
